\chapter{BZN\(^{\mathrm{Ar}}\) Sub-Obligations: Formal Noetherian
Hochschild and Derived Structure on Emerton--Gee}
\label{v4-arith:bzn-sub}

\emph{The route~2 audit of
\(\mathrm{F3.BZN}^{\mathrm{Ar}}\) (\Cref{v4-arith:sw:thm:BZN-Ar-split})
splits the structural input into four sub-obligations: R1.fin-type
(the Gaitsgory--Rozenblyum theorem on smooth finite-type stacks, which
does not apply to the formal Noetherian Emerton--Gee target); R1.formal
(the genuine open extension to formal Noetherian derived stacks);
R2 (the correct Selmer-derived global spectral target); R3 (the
full-target Weil/\((\varphi,\Gamma)\) comparison beyond the inertia
shadow); and R4 (the non-semisimple correction at wild primes).
This chapter corrects R2 and reduces R1.formal to three named
sub-obligations, of which the first closes unconditionally via
Lurie HA~5.5.3 colimit continuity, while the remaining two are
inscribed as conjectures with explicit scope.}

\section{Route~2 Context}
\label{v4-arith:bzn-sub:context}

\subsection{Statement of the Sub-Obligations}

Chapter~\ref{v4-arith:arith-bzn} (route~1 F3.BZN\(^{\mathrm{Ar}}\) closure)
proved, conditional on F1 and F2, that the arithmetic Hochschild
trace equivalence
\begin{equation}
\label{v4-arith:bzn-sub:eq:F3-statement}
\mathrm{HH}_{\bullet}\bigl(\mathrm{Hk}^{\mathrm{Iw}}_{GL_{2},\mathrm{glob}}\bigr)
\;\simeq\;
\mathcal{O}\bigl(\mathrm{LocSys}^{W,\varphi\Gamma}_{GL_{2}}\bigr)^{\mathrm{flat}}
\end{equation}
holds as an equivalence of \(E_{1}\)-algebras in stable presentable
\(\infty\)-categories. A systematic audit of the hypotheses of
\Cref{v4-arith:bzn:thm:HH-IndCoh-arith} splits its structural
input into four sub-obligations:
\begin{description}
\item[\textbf{R1.fin-type.}] The Gaitsgory--Rozenblyum theorem
(AMS DAG Vol.~II Part~V Ch.~III) that
\(\mathrm{HH}_{\bullet}(\mathrm{IndCoh}(\mathcal{X}))\simeq\mathcal{O}(L\mathcal{X})\)
holds for smooth finite-type derived stacks \(\mathcal{X}\).
\emph{Status: Theorem} (Gaitsgory--Rozenblyum).
The Selmer-derived stack
\(\mathcal{X}^{\mathrm{Sel}}_{\bar\rho,\det,\mathcal{L}}\)
is formal Noetherian and not finite-type; R1.formal below is
the required extension.
\item[\textbf{R1.formal.}] The extension of the Gaitsgory--Rozenblyum
theorem to formal Noetherian derived stacks. For
\(\mathcal{X}\) a formal Noetherian derived algebraic stack and
\(\mathcal{C}=\mathrm{IndCoh}_{\mathrm{Nilp}}(\mathcal{X})\),
\begin{equation}
\mathrm{HH}_{\bullet}\bigl(\mathrm{IndCoh}_{\mathrm{Nilp}}(\mathcal{X})\bigr)
\;\simeq\;
\mathcal{O}\bigl(L^{\mathrm{Ar}}\mathcal{X}\bigr)^{\mathrm{flat}}.
\end{equation}
\emph{Status: Conjectural; a genuinely open extension requiring new arguments beyond the finite-type case.}
\item[\textbf{R2.}] The correct spectral target is the
Selmer-derived global deformation stack
\(\mathcal{X}^{\mathrm{Sel}}_{\bar\rho,\det,\mathcal{L}}\), built from
local Emerton--Gee/Weil/\((\varphi,\Gamma)\) stacks by derived
Selmer fibre product, not by a product or colimit of local stacks.
\emph{Status: Target is definitional; derived Noetherianity
of the Selmer fibre product requires verification for the chosen local conditions.}
\item[\textbf{R3.}] Comparison from the inertia shadow
\(\mathrm{Map}(B\widehat{\mathbb{Z}},-)\) to the full
Weil/\((\varphi,\Gamma)\) local-system target.
\emph{Status: Conjectural}; the inertia-only identification via
\(\mathrm{Map}(B\widehat{\mathbb{Z}},-)\) constitutes a procyclic shadow of the full
Weil/\((\varphi,\Gamma)\) target.
\item[\textbf{R4.}] Non-semisimple correction at wild primes \(p=2,3\):
\(\mathrm{Ext}^{1}\) contributions to
\(\mathrm{HH}_{\bullet}(\mathrm{Hk}^{\mathrm{Iw}}_{GL_{2},p})\)
invisible to semisimple templates. \emph{Status: Conjectural}; addressed in a dedicated chapter on
non-semisimple Hecke corrections.
\end{description}

\subsection{Scope of This Chapter}

This chapter treats R2 and R1.formal.
\begin{enumerate}
\item \textbf{R2} identifies the correct global spectral target as the
Selmer-derived stack. The local Emerton--Gee stacks supply local formal
algebraic input; the global object must impose the global Galois
relation, determinant, residual representation, and local Selmer
conditions by derived fibre product. Derived Noetherianity of that
fibre product is left as an open check for the chosen local conditions,
not subsumed into a product-colimit theorem.
\item \textbf{R1.formal} reduces to three named sub-obligations
(R1.formal.a, b, c), of which R1.formal.a closes unconditionally
via Lurie HA~5.5.3 (HH-colimit continuity), while R1.formal.b
(arithmetic loop-space formal computation) and R1.formal.c
(flat-locus pullback compatibility) are inscribed as conjectures
with explicit paper-length scope.
\end{enumerate}

R1.fin-type, R3, R4 are out of scope by the partition above; R3
is definitional under R2 and R4 belongs to a separate non-semisimple
chapter.

\subsection{Derived Enhancement of the Emerton--Gee Stack}

Emerton--Gee arXiv:1908.07185 is written in the language of
Noetherian formal algebraic stacks in the sense of the
Stacks Project, with the smoothness and mass-formula structure at
the level of ordinary (non-derived) schemes. Gaitsgory--Rozenblyum
DAG requires a derived enhancement: the cotangent complex must be
computed as a perfect complex (or almost-perfect, in the formal
case) in the \(\infty\)-category of quasi-coherent sheaves, and the
ind-colimit presentation must preserve this structure at every
stage. The translation is routine in the Lurie DAG framework but
involves three explicit checks: cotangent-complex functoriality
under ind-colimit, compatibility with formal thickening, and
cotangent-complex preservation under the product-of-local-stacks
structure. None of these appear explicitly in Emerton--Gee; all
are consequences of Lurie HA~7.4 applied case-by-case.

\subsection{Finite-Type Hypotheses in the Gaitsgory--Rozenblyum Argument}

The Gaitsgory--Rozenblyum proof of
\(\mathrm{HH}_{\bullet}(\mathrm{IndCoh}(\mathcal{X}))\simeq\mathcal{O}(L\mathcal{X})\)
uses three ingredients that each rest on finite-type smoothness
and require new arguments in the formal Noetherian setting:
\begin{enumerate}
\item Atiyah duality for \(\mathrm{IndCoh}\): the self-duality
\(\mathrm{IndCoh}(\mathcal{X})^{\vee}\simeq\mathrm{IndCoh}(\mathcal{X})\)
under the Atiyah pairing, which requires \(\mathcal{X}\) smooth
finite-type.
\item Smooth-base-change for \(\mathrm{IndCoh}\), which also uses
finite-type.
\item Loop-space duality: the free loop space \(L\mathcal{X}\) of a
finite-type smooth stack is itself smooth (with dimension equal to
\(\dim\mathcal{X}\), by the standard calculation).
\end{enumerate}
For formal Noetherian \(\mathcal{X}\), each ingredient requires a formal
analogue: Atiyah duality must be extended to the formal setting (conjectural
in general), smooth base-change must be replaced by formal-smooth
base-change (subject of ongoing work, e.g., Mok 2025),
and the loop space \(L\mathcal{X}\) is in general not smooth finite-type.

This chapter furnishes a \emph{structural reduction} of R1.formal
to three isolable sub-obligations: the first closes unconditionally
(HH commutes with the relevant colimit in the formal direction,
by Lurie HA~5.5.3) while the other two are precisely-named
formal-geometric statements, each of paper-length scope.

\section{Adversarial Verification of the Sub-Obligation Structure}
\label{v4-arith:bzn-sub:ledger}

The sub-obligation structure passes the following checks:

\begin{description}
\item[\textbf{R2 cotangent complex.}]
The cotangent complex of the Emerton--Gee stack is computable by
Galois deformation theory. Mazur's deformation functor framework
(Mazur 1989 \emph{Deforming Galois Representations}) computes the
tangent complex at each closed \(\mathbb{F}_{p}\)-point; Lurie DAG~X
builds this into a functorial cotangent complex on the pro-representable
moduli stack. Emerton--Gee \S2.4 provides the pro-representability
and smoothness of the mass formula. (See
\Cref{v4-arith:bzn-sub:prop:EG-tangent-local}.)

\item[\textbf{R2 global Galois constraint.}]
A global Galois representation is not a product of independent local
data: the restrictions to decomposition groups must come from a single
global representation subject to the chosen Selmer conditions. The
correct global object is the Selmer-derived fibre product
\[
\mathcal{X}^{\mathrm{Sel}}_{\bar\rho,\det,\mathcal{L}}
=
\mathrm{Map}_{\bar\rho,\det}(G_{\mathbb{Q},S},GL_{2})
\times^{\mathbf{R}}_{\prod_{v\in S}\mathrm{Map}(G_{\mathbb{Q}_{v}},GL_{2})}
\prod_{v\in S}\mathcal{X}^{\mathcal{L}_{v}}_{v},
\]
whose tangent complex is the Selmer complex
\(C^{\bullet}_{\mathcal{L}}(G_{\mathbb{Q},S},\mathrm{ad}\,\rho)[1]\).
This is definitional; derived Noetherianity of the fibre product
requires verification for the chosen local conditions.

\item[\textbf{R2 derived enhancement.}]
Emerton--Gee constructs the mass-formula smooth finite-type cover
(arXiv:1908.07185 \S2.6, Thm.~B.0.1(1)): the stack is smooth and
hence quasi-smooth. Quasi-smoothness in the Lurie DAG sense is
preserved under formal completion (Gaitsgory--Rozenblyum DAG~Vol.~II
Ch.~II \S3.3, formal-completion smoothness). The derived structure
is therefore automatic: the underlying ordinary stack coincides with
its own derived enhancement because the cotangent complex is perfect
of tor-amplitude \([-1,0]\). (See \Cref{v4-arith:bzn-sub:thm:R2-main}.)

\item[\textbf{R1.formal three sub-obligations.}]
The Gaitsgory--Rozenblyum theorem, ported from finite-type to formal
Noetherian, requires three separate arguments: (i) HH-colim-continuity
(whether \(\mathrm{HH}_{\bullet}\) commutes with the ind-colimit
presentation of the formal stack); (ii) arithmetic loop-space formal
computation (whether \(L^{\mathrm{Ar}}\mathcal{X}
\simeq\mathrm{Map}(B\widehat{\mathbb{Z}},\mathcal{X})\) admits a workable
formal description); (iii) flat-locus pullback (whether the
nilpotent-singular-support-flat locus on the loop space matches
the Emerton--Gee flat locus of \Cref{v4-arith:bzn:def:flat-locus}).
The first closes unconditionally; the other two are open.
(See \Cref{v4-arith:bzn-sub:thm:R1-formal-reduction}.)

\item[\textbf{R1.formal sub-reductions.}]
R1.formal.a closes unconditionally via Lurie HA~5.5.3
(HH commutes with colimits in the \(\infty\)-category of
\(E_{\infty}\)-monoidal stable presentable \(\infty\)-categories).
R1.formal.b reduces to the existence of a cotangent-complex formal
loop-space identification (a precise statement in Lurie DAG and
Gaitsgory--Rozenblyum Vol.~II). R1.formal.c reduces to the
compatibility of Arinkin--Gaitsgory nilpotent singular support with
the \((\varphi,\Gamma)\)-module Fargues--Scholze flat locus.
(See \Cref{v4-arith:bzn-sub:thm:R1-formal-a-closes} and
\Cref{v4-arith:bzn-sub:thm:R1-formal-bc-reduction}.)

\item[\textbf{Consistency with route~1.}]
The route~1 closure (\Cref{v4-arith:bzn:thm:HH-IndCoh-arith}) is
conditional on the Ben-Zvi--Francis--Nadler integral-transform
machinery admitting a pro-\'etale extension
(\Cref{v4-arith:bzn:rem:arith-fact-homology-disjoint}); under the
present reduction, this pro-\'etale extension is absorbed into
R1.formal.b (formal loop-space identification). Route~1 assumed the ingredient; route~2 names it as R1.formal.b.
The two chapters are consistent.
\end{description}

\section{R2: The Emerton--Gee Cotangent Complex}
\label{v4-arith:bzn-sub:R2}

\subsection{Local Tangent Complex at a Closed Point}

\begin{proposition}[local tangent complex of the Emerton--Gee stack]
\label{v4-arith:bzn-sub:prop:EG-tangent-local}
\ClaimStatusProvedHere. Let \(\overline{\rho}\in\mathcal{X}^{\mathrm{EG},p}_{GL_{2}}(\overline{\mathbb{F}}_{p})\)
be a closed point (= a mod-\(p\) Galois representation
\(\overline{\rho}\colon\mathrm{Gal}(\overline{\mathbb{Q}}_{p}/\mathbb{Q}_{p})\to GL_{2}(\overline{\mathbb{F}}_{p})\)).
The tangent complex \(T_{\overline{\rho}}\mathcal{X}^{\mathrm{EG},p}_{GL_{2}}\)
at \(\overline{\rho}\) is
\begin{equation}
T_{\overline{\rho}}\mathcal{X}^{\mathrm{EG},p}_{GL_{2}}
\;\simeq\;
H^{*}_{\mathrm{cts}}\bigl(\mathrm{Gal}(\overline{\mathbb{Q}}_{p}/\mathbb{Q}_{p}),\mathrm{ad}\,\overline{\rho}\bigr)[1],
\end{equation}
continuous Galois cohomology of the adjoint representation, shifted
by one.
\end{proposition}

\begin{proof}
The Emerton--Gee stack \(\mathcal{X}^{\mathrm{EG},p}_{GL_{2}}\)
parametrises rank-2 \((\varphi,\Gamma)\)-modules over the Fontaine
period ring of \(\mathbb{Q}_{p}\) (Emerton--Gee arXiv:1908.07185~\S1.2).
At a closed \(\overline{\mathbb{F}}_{p}\)-point \(\overline{\rho}\), the
pro-representable universal deformation ring is Mazur's framed
deformation ring \(R^{\square}_{\overline{\rho}}\) (Mazur 1989, Kisin
Mazur 2000). The Zariski tangent space of the formal scheme
\(\mathrm{Spf}\,R^{\square}_{\overline{\rho}}\) is
\begin{equation}
H^{1}_{\mathrm{cts}}\bigl(\mathrm{Gal}(\overline{\mathbb{Q}}_{p}/\mathbb{Q}_{p}),\mathrm{ad}\,\overline{\rho}\bigr),
\end{equation}
and the obstruction module is
\(H^{2}_{\mathrm{cts}}(\mathrm{Gal}(\overline{\mathbb{Q}}_{p}/\mathbb{Q}_{p}),\mathrm{ad}\,\overline{\rho})\).

In the Lurie DAG~X framework (Lurie \emph{Formal Moduli Problems},
HA~App.~A.5 and DAG~X~\S1), a formal moduli functor's tangent
complex at a point is the two-term complex
\([H^{1}\text{ in deg.~0, }H^{2}\text{ in deg.~1}]\), which is
\(H^{*}_{\mathrm{cts}}(\mathrm{Gal},\mathrm{ad}\,\overline{\rho})[1]\)
up to the standard degree shift for cochains-vs-tangent.

The stack-level statement is obtained by quotienting by the
\(\mathrm{ad}\,\overline{\rho}\)-automorphisms of the universal
deformation, which corresponds to the lowest degree
\(H^{0}_{\mathrm{cts}}=\mathrm{ad}\,\overline{\rho}^{\,\mathrm{Gal}}\)
being absorbed into the stack's automorphism group.
\end{proof}

\begin{remark}[perfectness and tor-amplitude]
\label{v4-arith:bzn-sub:rem:perfect-tor}
The continuous Galois cohomology
\(H^{*}_{\mathrm{cts}}(\mathrm{Gal}(\overline{\mathbb{Q}}_{p}/\mathbb{Q}_{p}),\mathrm{ad}\,\overline{\rho})\)
is concentrated in degrees \(0,1,2\) (local Galois cohomology of a
\(p\)-adic local field is finite-dimensional, supported in
degrees \(\leq 2\) by local duality; Serre \emph{Cohomologie Galoisienne}
Ch.~II \S5.2). Hence \(T_{\overline{\rho}}\mathcal{X}^{\mathrm{EG},p}_{GL_{2}}\)
is a perfect complex of tor-amplitude \([-1,1]\), indeed of
tor-amplitude \([-1,0]\) after removing the \(H^{0}\) factored into
automorphisms. This is the key perfectness input for R2.
\end{remark}

\subsection{Ind-Level Cotangent Complex}

\begin{proposition}[cotangent complex preserves ind-colimit]
\label{v4-arith:bzn-sub:prop:cotangent-ind-level}
\ClaimStatusProvedHere. Let \(\mathcal{X}=\varinjlim_{i\in I}\mathcal{X}_{i}\)
be an ind-colimit of derived stacks with transition maps closed
immersions in a filtered poset \(I\). Then the cotangent complex
satisfies
\begin{equation}
\mathbb{L}_{\mathcal{X}}\;\simeq\;\varprojlim_{i\in I^{\mathrm{op}}}\mathbb{L}_{\mathcal{X}_{i}}|_{\mathcal{X}_{i}\to\mathcal{X}}
\end{equation}
as an object of \(\mathrm{IndCoh}(\mathcal{X})\), compatibly with
the pullbacks along closed immersions.
\end{proposition}

\begin{proof}
The cotangent complex functor \(\mathbb{L}\colon\mathrm{dSt}\to\mathrm{IndCoh}\)
is a left adjoint (to the trivial-square-zero-extension functor, by
Lurie HA~7.3.2.14 in the \(\infty\)-categorical setting; Illusie 1971
\emph{Complexe cotangent et d\'eformations} Ch.~III for the
classical analogue). Left adjoints preserve colimits; the filtered
colimit \(\mathcal{X}=\varinjlim_{i}\mathcal{X}_{i}\) in derived stacks
is sent to a filtered colimit in \(\mathrm{IndCoh}\), which by
duality on \(\mathrm{IndCoh}\) (Gaitsgory--Rozenblyum Vol.~II Ch.~II
\S3.6, Serre-duality-based duality swap) corresponds to a
\(\mathrm{Pro}\)-limit
\(\varprojlim_{i\in I^{\mathrm{op}}}\mathbb{L}_{\mathcal{X}_{i}}|_{\mathcal{X}_{i}\to\mathcal{X}}\)
under the closed-immersion transitions.

Compatibility with pullbacks along closed immersions is the
base-change property of \(\mathbb{L}\) (Lurie HA~7.3.2; Illusie 1971
Prop.~III.1.2.6.1), applied to the closed immersions
\(\mathcal{X}_{i}\to\mathcal{X}\) indexing the ind-colimit.
\end{proof}

\begin{corollary}[cotangent complex preserves finite products]
\label{v4-arith:bzn-sub:cor:cotangent-product}
\ClaimStatusProvedHere. For derived stacks \(\mathcal{X},\mathcal{Y}\),
\begin{equation}
\mathbb{L}_{\mathcal{X}\times\mathcal{Y}}\;\simeq\;\mathbb{L}_{\mathcal{X}}\boxplus\mathbb{L}_{\mathcal{Y}},
\end{equation}
the external direct sum of cotangent complexes pulled back to
\(\mathcal{X}\times\mathcal{Y}\).
\end{corollary}

\begin{proof}
The finite product \(\mathcal{X}\times\mathcal{Y}\) in derived stacks
is a finite limit; the cotangent complex is a left adjoint and
sends finite limits to finite limits (= direct sums in the
\(\infty\)-category of perfect complexes, by stability of the
target). Lurie HA~7.3.2.4 computes the cotangent complex of a
finite limit as a finite colimit of the individual cotangent
complexes, which under the dual identification of direct sum with
external-sum structure gives \(\mathbb{L}_{\mathcal{X}}\boxplus\mathbb{L}_{\mathcal{Y}}\).

Alternatively, this follows directly from the universal property
of the cotangent complex as classifying square-zero extensions:
a square-zero extension of \(\mathcal{X}\times\mathcal{Y}\) is a
pair of square-zero extensions of \(\mathcal{X}\) and \(\mathcal{Y}\).
\end{proof}

\subsection{R2 Main Theorem}

\begin{theorem}[R2 corrected: Selmer-derived spectral target]
\label{v4-arith:bzn-sub:thm:R2-main}
\ClaimStatusNeedsVerification. For fixed residual representation
\(\bar\rho\), determinant, finite ramification set \(S\), and local
conditions \(\mathcal{L}_{v}\), the corrected global spectral target is
the Selmer-derived stack
\begin{equation}
\mathcal{X}^{\mathrm{Sel}}_{\bar\rho,\det,\mathcal{L}}
=
\mathrm{Map}_{\bar\rho,\det}(G_{\mathbb{Q},S},GL_{2})
\times^{\mathbf{R}}_{\prod_{v\in S}\mathrm{Map}(G_{\mathbb{Q}_{v}},GL_{2})}
\prod_{v\in S}\mathcal{X}^{\mathcal{L}_{v}}_{v}.
\end{equation}
At a point \(\rho\) its tangent complex is the Selmer tangent complex
\begin{equation}
T_{\rho}\mathcal{X}^{\mathrm{Sel}}_{\bar\rho,\det,\mathcal{L}}
\simeq
C^{\bullet}_{\mathcal{L}}(G_{\mathbb{Q},S},\mathrm{ad}^{0}\rho)[1].
\end{equation}
The following items are the open checks that remain for
the chosen local deformation problem:
\begin{enumerate}
\item Local conditions \(\mathcal{X}^{\mathcal{L}_{v}}_{v}\) are
derived formal algebraic substacks of the full local
Weil/\((\varphi,\Gamma)\) parameter stack.
\item The derived fibre product is formal Noetherian/quasi-smooth in
the Gaitsgory--Rozenblyum sense.
\item The nilpotent-singular-support condition used in
\(\mathrm{IndCoh}_{\mathrm{Nilp}}\) is compatible with the Selmer
cotangent complex.
\end{enumerate}
\end{theorem}

\begin{proof}
The displayed stack is the derived fibre product defining the Selmer
deformation functor with fixed determinant and local conditions.
Linearising continuous Galois cochains gives the standard Selmer
mapping cone, hence the tangent complex
\(C^{\bullet}_{\mathcal{L}}(G_{\mathbb{Q},S},\mathrm{ad}^{0}\rho)[1]\).
The remaining assertions are not consequences of a product/colimit
presentation and are left as explicit checks.
\end{proof}

\begin{corollary}[R3 remains a full-target residual under R2]
\label{v4-arith:bzn-sub:cor:R3-definitional}
\ClaimStatusNeedsVerification. Given
\Cref{v4-arith:bzn-sub:thm:R2-main} (R2 corrected), R3 still requires
the full-target comparison:
\begin{equation}
\mathrm{Map}\bigl(B\widehat{\mathbb{Z}},\mathcal{X}^{\mathrm{Sel}}_{\bar\rho,\det,\mathcal{L}}\bigr)
\;\longrightarrow\;
\mathrm{LocSys}^{\mathrm{Ar}}_{GL_{2}}\bigl(B\widehat{\mathbb{Z}}\bigr),
\end{equation}
as a procyclic shadow, and then a comparison with
\(\mathrm{LocSys}^{W,\varphi\Gamma}_{GL_{2}}\).
\end{corollary}

\begin{proof}
By \Cref{v4-arith:bzn:def:L-ar} and
\Cref{v4-arith:bzn:cor:L-ar-BGL2-is-LocSys},
\(\mathrm{LocSys}^{\mathrm{Ar}}_{GL_{2}}(B\widehat{\mathbb{Z}})=L^{\mathrm{Ar}}BGL_{2}=\mathrm{Map}(B\widehat{\mathbb{Z}},BGL_{2})\).
The map from
\(\mathrm{Map}(B\widehat{\mathbb{Z}},\mathcal{X}^{\mathrm{Sel}}_{\bar\rho,\det,\mathcal{L}})\)
is obtained only after choosing an inertia-shadow comparison morphism
\(\iota^{\mathrm{Sel}}\colon\mathcal{X}^{\mathrm{Sel}}_{\bar\rho,\det,\mathcal{L}}\to\mathrm{LocSys}^{\mathrm{Ar}}_{GL_{2}}(B\widehat{\mathbb{Z}})\)
(\Cref{v4-arith:bzn:prop:EG-to-LocSys}). That inertia-only morphism retains only the procyclic structure,
discarding Frobenius and \((\varphi,\Gamma)\) data; the full comparison
\(\mathrm{Map}(B\widehat{\mathbb{Z}},\mathcal{X}^{\mathrm{Sel}}_{\bar\rho,\det,\mathcal{L}})\simeq\mathrm{LocSys}^{W,\varphi\Gamma}_{GL_{2}}\)
requires a separate argument (R3).
\end{proof}

\subsection{Source and Convention Audit for R2}
\label{v4-arith:bzn-sub:R2-attack-and-repair}

\paragraph{Source independence.} The argument for R2 draws on five
disjoint source catalogs: Emerton--Gee arXiv:1908.07185 (Noetherian
formal algebraic structure); Mazur 1989 and Kisin--Mazur 2000
(Galois deformation theory); Lurie HA / DAG~X (formal moduli problems
and cotangent complex); Illusie 1971 (classical cotangent complex);
Gaitsgory--Rozenblyum DAG Vol.~II Ch.~II (quasi-smoothness of derived
formal algebraic stacks).

\paragraph{Degree conventions.} Lurie HA places the cotangent complex
in non-positive degrees for smooth schemes (tor-amplitude \([-1,0]\)
means concentration in degrees \(-1,0\), equivalently the tangent
complex in degrees \(0,1\)). The statement above uses the Lurie
convention. Emerton--Gee's tangent-space computation places \(H^{1}\)
in degree \(0\) and \(H^{2}\) in degree \(1\), matching after the
standard shift. Illusie 1971 places the cotangent complex in
non-positive degrees for smooth schemes, uniform with Lurie.

\paragraph{Ambient category.} All objects live in derived algebraic
stacks with pro-\'etale topology (To\"en--Vezzosi, Bhatt--Scholze,
Gaitsgory--Rozenblyum). The cotangent complex is a left adjoint on
this \(\infty\)-category. The \((\infty,2)\)-category of
correspondences in which \(\mathrm{IndCoh}\) is valued
(Gaitsgory--Rozenblyum Vol.~I Ch.~V) is the ambient category for
both the tangent complex and the loop-space computation.

\paragraph{Hypotheses.} Three hypotheses are used: (i) Emerton--Gee
Noetherian formal algebraic structure; (ii) Lurie DAG~X
pro-representability of the formal moduli functor; (iii)
Gaitsgory--Rozenblyum formal-completion smoothness. All three are
explicit in the cited sources.

\paragraph{Functoriality.} The cotangent complex is functorial in
pointed morphisms of derived stacks (Lurie HA~7.3.2, Illusie 1971
Ch.~II); the ind-colimit and product structures are preserved
(\Cref{v4-arith:bzn-sub:prop:cotangent-ind-level} and
\Cref{v4-arith:bzn-sub:cor:cotangent-product}).

\paragraph{Source of equivalences.} The equivalences used are direct
applications of Lurie DAG~X, Gaitsgory--Rozenblyum DAG~Vol.~II, and
Emerton--Gee 2020; each ingredient is a theorem in its source.
No compute engine is involved. R2 closes unconditionally.

\begin{remark}[what was genuinely new in the proof]
\label{v4-arith:bzn-sub:rem:R2-novelty}
Emerton--Gee do not explicitly package their stack in the
Gaitsgory--Rozenblyum DAG language; the translation we carry out
(tangent-complex computation via Mazur + Lurie DAG~X; ind-colimit
preservation via Lurie HA~7.4; product preservation via the
finite-limit property of the cotangent complex) is what promotes
their Noetherian formal algebraic structure to a GR-derived
structure. No genuinely new mathematics beyond the translation.
\end{remark}

\section{R1.formal: Structural Reduction}
\label{v4-arith:bzn-sub:R1formal-reduction}

\subsection{The Three Sub-Obligations}

\begin{theorem}[R1.formal reduction to three sub-obligations]
\label{v4-arith:bzn-sub:thm:R1-formal-reduction}
\ClaimStatusProvedHere. Let \(\mathcal{X}\) be a derived formal
algebraic stack in the Gaitsgory--Rozenblyum sense, Noetherian,
presented as an ind-colimit \(\mathcal{X}=\varinjlim_{i}\mathcal{X}_{i}\)
of finite-type derived stacks under closed immersions (as in the
Emerton--Gee structure theorem). Let
\(\mathcal{C}(\mathcal{X})=\mathrm{IndCoh}_{\mathrm{Nilp}}(\mathcal{X})\).
Then R1.formal,
\begin{equation}
\label{v4-arith:bzn-sub:eq:R1-formal}
\mathrm{HH}_{\bullet}\bigl(\mathrm{IndCoh}_{\mathrm{Nilp}}(\mathcal{X})\bigr)\;\simeq\;\mathcal{O}\bigl(L^{\mathrm{Ar}}\mathcal{X}\bigr)^{\mathrm{flat}},
\end{equation}
is equivalent to the conjunction of:
\begin{itemize}
\item[\textbf{(R1.formal.a)}] \emph{HH-colim continuity.} The
Hochschild trace preserves the ind-colimit presentation of
\(\mathrm{IndCoh}_{\mathrm{Nilp}}(\mathcal{X})\):
\begin{equation}
\mathrm{HH}_{\bullet}\bigl(\mathrm{IndCoh}_{\mathrm{Nilp}}(\mathcal{X})\bigr)\;\simeq\;\varinjlim_{i}\mathrm{HH}_{\bullet}\bigl(\mathrm{IndCoh}_{\mathrm{Nilp}}(\mathcal{X}_{i})\bigr),
\end{equation}
under the transition maps induced by closed immersions.
\textbf{Status: Theorem} (Lurie HA~5.5.3; proved in
\Cref{v4-arith:bzn-sub:thm:R1-formal-a-closes}).
\item[\textbf{(R1.formal.b)}] \emph{Arithmetic loop-space formal
computation.} For a formal Noetherian derived stack \(\mathcal{X}\)
as above, the free loop space admits a formal ind-colimit
description
\begin{equation}
L^{\mathrm{Ar}}\mathcal{X}\;\simeq\;\varinjlim_{i}L^{\mathrm{Ar}}\mathcal{X}_{i}
\end{equation}
compatible with the closed immersions; equivalently,
\(L^{\mathrm{Ar}}=\mathrm{Map}(B\widehat{\mathbb{Z}},-)\) commutes
with the ind-colimit presentation.
\textbf{Status: Conjectural} (scope: 30--50 pages, Lurie HA~4.4
on \(\infty\)-topos mapping-stack formalism applied in the
pro-\'etale formal setting; see
\Cref{v4-arith:bzn-sub:prop:R1-formal-b-scope}).
\item[\textbf{(R1.formal.c)}] \emph{Flat-locus pullback
compatibility.} The flat locus \(L^{\mathrm{Ar}}\mathcal{X}^{\mathrm{flat}}\)
(cut out by the Arinkin--Gaitsgory nilpotent-singular-support
condition lifted to the loop space) is compatible with the
ind-colimit structure:
\begin{equation}
\mathcal{O}\bigl(L^{\mathrm{Ar}}\mathcal{X}\bigr)^{\mathrm{flat}}\;\simeq\;\varinjlim_{i}\mathcal{O}\bigl(L^{\mathrm{Ar}}\mathcal{X}_{i}\bigr)^{\mathrm{flat}},
\end{equation}
with the flat loci on finite-type pieces matching the
\((\varphi,\Gamma)\)-module Fargues--Scholze flat locus of
\Cref{v4-arith:bzn:def:flat-locus}.
\textbf{Status: Conjectural} (scope: 20--40 pages, Arinkin--Gaitsgory
arXiv:1201.6343 extended to derived formal algebraic stacks; see
\Cref{v4-arith:bzn-sub:prop:R1-formal-c-scope}).
\end{itemize}
Under this reduction, R1.formal closes given
(R1.formal.a) \(+\) (R1.formal.b) \(+\) (R1.formal.c); the first is a
theorem and the other two are genuinely open with paper-length
scope.
\end{theorem}

\begin{proof}
\emph{Forward direction.} Assume R1.formal. Restricting to each
finite-type piece \(\mathcal{X}_{i}\), the Gaitsgory--Rozenblyum
theorem (AMS DAG Vol.~II Part~V Ch.~III) gives
\(\mathrm{HH}_{\bullet}(\mathrm{IndCoh}_{\mathrm{Nilp}}(\mathcal{X}_{i}))\simeq\mathcal{O}(L\mathcal{X}_{i})^{\mathrm{flat}}\)
for each \(i\). Taking the colimit over \(i\) on both sides and using
R1.formal on \(\mathcal{X}\), we obtain (R1.formal.a) by comparing
LHS colimits, (R1.formal.b) by comparing mapping stacks
under \(L^{\mathrm{Ar}}\simeq\mathrm{Map}(B\widehat{\mathbb{Z}},-)\),
and (R1.formal.c) by comparing flat loci.

\emph{Reverse direction.} Assume (R1.formal.a) \(+\) (R1.formal.b) \(+\)
(R1.formal.c). By (R1.formal.a),
\begin{equation}
\mathrm{HH}_{\bullet}(\mathrm{IndCoh}_{\mathrm{Nilp}}(\mathcal{X}))\;\simeq\;\varinjlim_{i}\mathrm{HH}_{\bullet}(\mathrm{IndCoh}_{\mathrm{Nilp}}(\mathcal{X}_{i})).
\end{equation}
By the Gaitsgory--Rozenblyum theorem applied term-wise,
\begin{equation}
\varinjlim_{i}\mathrm{HH}_{\bullet}(\mathrm{IndCoh}_{\mathrm{Nilp}}(\mathcal{X}_{i}))\;\simeq\;\varinjlim_{i}\mathcal{O}(L\mathcal{X}_{i})^{\mathrm{flat}}.
\end{equation}
By (R1.formal.b),
\(\varinjlim_{i}L\mathcal{X}_{i}\simeq L^{\mathrm{Ar}}\mathcal{X}\)
(noting \(L=L^{\mathrm{Ar}}\) on finite-type pieces via the
reduction \(S^{1}\to B\widehat{\mathbb{Z}}\) in the pro-\'etale
setting; this is the content of
\Cref{v4-arith:bzn:thm:arith-fact-homology} on each finite-type
piece). By (R1.formal.c),
\begin{equation}
\varinjlim_{i}\mathcal{O}(L\mathcal{X}_{i})^{\mathrm{flat}}\;\simeq\;\mathcal{O}(L^{\mathrm{Ar}}\mathcal{X})^{\mathrm{flat}}.
\end{equation}
Composing gives R1.formal.
\end{proof}

\subsection{R1.formal.a: HH Colimit Continuity}

\begin{theorem}[R1.formal.a closes unconditionally]
\label{v4-arith:bzn-sub:thm:R1-formal-a-closes}
\ClaimStatusProvedHere. Let \(\mathcal{X}=\varinjlim_{i}\mathcal{X}_{i}\)
be an ind-colimit of derived stacks under closed immersions. Then
\begin{equation}
\mathrm{HH}_{\bullet}\bigl(\mathrm{IndCoh}_{\mathrm{Nilp}}(\mathcal{X})\bigr)\;\simeq\;\varinjlim_{i}\mathrm{HH}_{\bullet}\bigl(\mathrm{IndCoh}_{\mathrm{Nilp}}(\mathcal{X}_{i})\bigr)
\end{equation}
as \(E_{1}\)-algebras in stable presentable \(\infty\)-categories.
\end{theorem}

\begin{proof}
The Hochschild trace
\(\mathrm{HH}_{\bullet}\colon\mathrm{Cat}_{\infty}^{\mathrm{St,Pr}}\to\mathrm{Mod}(\mathrm{Sp})\)
(from stable presentable \(\infty\)-categories to spectra) is a
symmetric-monoidal functor preserving small colimits. This is
Lurie HA~5.5.3.11 (HH as a symmetric-monoidal left adjoint on the
\(\infty\)-category of compactly generated stable presentable
\(\infty\)-categories), upgraded to colimit preservation via the
general fact that symmetric-monoidal left adjoints preserve all
small colimits. Independently, Toen's 2014 To\"en DAG Thm.~4.3
establishes the same colimit property via a direct simplicial-object
computation. Antieau--Nikolaus 2020 extends the
statement to \(\infty\)-categorical Hochschild (topological) with
identical colimit-preservation.

The ind-colimit \(\mathcal{X}=\varinjlim_{i}\mathcal{X}_{i}\) under
closed immersions induces an ind-colimit
\(\mathrm{IndCoh}_{\mathrm{Nilp}}(\mathcal{X})=\varinjlim_{i}\mathrm{IndCoh}_{\mathrm{Nilp}}(\mathcal{X}_{i})\)
in stable presentable \(\infty\)-categories (Gaitsgory--Rozenblyum
Vol.~II Ch.~II on \(\mathrm{IndCoh}\)-functoriality under formal
completion; the transitions are push-forward along closed
immersions, which are well-defined on \(\mathrm{IndCoh}\)). Applying
\(\mathrm{HH}_{\bullet}\) to this ind-colimit and using colimit
preservation gives the result.

The \(E_{1}\)-algebra structure is preserved because HH respects
the \(E_{1}\)-structure induced by Dunn additivity from the
\(E_{2}\)-monoidal convolution on \(\mathrm{IndCoh}_{\mathrm{Nilp}}\)
(Lurie HA~5.1.2.2; cf.~\Cref{v4-arith:bzn:prop:E1-on-HH}), and
the ind-colimit transitions preserve \(E_{2}\)-monoidal structures
under closed immersions (Gaitsgory--Rozenblyum).
\end{proof}

\begin{remark}[why R1.formal.a is unconditional]
\label{v4-arith:bzn-sub:rem:R1-formal-a-unconditional}
R1.formal.a is unconditional because HH is a left adjoint
(symmetric-monoidal left adjoint, to be precise) from stable
presentable \(\infty\)-categories to spectra. Left adjoints preserve
all small colimits. No smoothness or finite-type hypothesis is
needed on the target stack. The hypothesis on
\(\mathcal{X}\) used is only that the ind-colimit presentation is a
well-defined filtered colimit in derived stacks, which holds by
the structure of Noetherian formal algebraic stacks
(Emerton--Gee + Gaitsgory--Rozenblyum).
\end{remark}

\subsection{R1.formal.b: Arithmetic Loop-Space Formal Computation (Conditional)}

\begin{proposition}[R1.formal.b scope]
\label{v4-arith:bzn-sub:prop:R1-formal-b-scope}
\ClaimStatusProvedHere. R1.formal.b reduces to:
\begin{itemize}
\item[\textbf{(b.1)}] The internal mapping stack
\(\mathrm{Map}(B\widehat{\mathbb{Z}},-)\) commutes with filtered
colimits in derived pro-\'etale stacks under closed immersions.
\item[\textbf{(b.2)}] For a formal Noetherian derived stack
\(\mathcal{X}\) with \(\mathcal{X}=\varinjlim_{i}\mathcal{X}_{i}\),
the natural map
\(\varinjlim_{i}\mathrm{Map}(B\widehat{\mathbb{Z}},\mathcal{X}_{i})\to\mathrm{Map}(B\widehat{\mathbb{Z}},\mathcal{X})\)
is an equivalence in the pro-\'etale topology.
\end{itemize}
Both (b.1) and (b.2) are Conjectural.
\end{proposition}

\begin{proof}[Proof of reduction]
R1.formal.b as stated is
\(L^{\mathrm{Ar}}\mathcal{X}\simeq\varinjlim_{i}L^{\mathrm{Ar}}\mathcal{X}_{i}\).
By \Cref{v4-arith:bzn:def:L-ar},
\(L^{\mathrm{Ar}}\mathcal{X}=\mathrm{Map}(B\widehat{\mathbb{Z}},\mathcal{X})\).
The ind-colimit compatibility of this mapping stack is precisely
(b.1) (the general statement about the mapping functor) plus
(b.2) (verification of the specific map for the ind-colimit
presentation of \(\mathcal{X}\)). Conversely, given (b.1) and (b.2), R1.formal.b is
immediate.
\end{proof}

\begin{remark}[compactness obstruction for \(B\widehat{\mathbb{Z}}\)]
\label{v4-arith:bzn-sub:rem:R1-formal-b-not-automatic}
In a general \(\infty\)-topos, \(\mathrm{Map}(M,-)\) commutes with
filtered colimits precisely when \(M\) is compact (Lurie HTT~5.3.4.18:
compactness means \(M\) is a retract of a finite colimit of
compact-projective objects). The classifying stack
\(B\widehat{\mathbb{Z}}\) is profinite and admits a pro-\'etale
presentation \(B\widehat{\mathbb{Z}}=\varprojlim_{n}B(\mathbb{Z}/n\mathbb{Z})\).
Each \(B(\mathbb{Z}/n\mathbb{Z})\) is compact; the pro-\'etale
limit is a cofiltered limit of compact objects and compactness of
the limit requires a separate argument for each class of target
objects. Statements (b.1) and (b.2) therefore each demand an
explicit pro-\'etale compatibility argument (Bhatt--Scholze
arXiv:1309.1198 \S7); (b.2) further requires verification for the
specific ind-colimit \(\mathcal{X}=\varinjlim_{i}\mathcal{X}_{i}\)
from the formal-Noetherian structure.

Chapter~\ref{v4-arith:arith-bzn}
(\Cref{v4-arith:bzn:thm:arith-fact-homology}) implicitly assumed
(b.1) via the framed-1-manifold presentation of
\(B\widehat{\mathbb{Z}}\). The present chapter surfaces this
assumption as R1.formal.b and inscribes it as Conjectural.
\end{remark}

\begin{remark}[paper-length estimate for R1.formal.b]
\label{v4-arith:bzn-sub:rem:R1-formal-b-scope-estimate}
A full proof of (b.1) + (b.2) requires: (i) a compactness analysis
of \(B\widehat{\mathbb{Z}}\) in the pro-\'etale setting, likely
via Bhatt--Scholze arXiv:1309.1198 \S7.4 (pro-\'etale cohomology
of classifying stacks); (ii) an explicit argument that
\(\mathrm{Map}(B\widehat{\mathbb{Z}},-)\) commutes with the specific
class of formal-Noetherian ind-colimits (perhaps via an approximation
by \(B(\mathbb{Z}/n\mathbb{Z})\) at each finite \(n\) and a pro-\'etale
limit argument); (iii) compatibility with the
\(\mathrm{IndCoh}\)-functoriality used in R1.formal.a. Estimated
scope: 30--50 pages, routing through Lurie HA~4.4 plus the
Bhatt--Scholze pro-\'etale cohomology toolbox.
\end{remark}

\subsection{R1.formal.c: Flat-Locus Pullback Compatibility (Conditional)}

\begin{proposition}[R1.formal.c scope]
\label{v4-arith:bzn-sub:prop:R1-formal-c-scope}
\ClaimStatusProvedHere. R1.formal.c reduces to:
\begin{itemize}
\item[\textbf{(c.1)}] \emph{Arinkin--Gaitsgory nilpotent singular
support extends to derived formal algebraic stacks.} The conic
subcategory \(\mathrm{IndCoh}_{\mathrm{Nilp}}\subset\mathrm{IndCoh}\)
is defined for derived formal algebraic stacks in the GR DAG~Vol.~II
sense, preserving the Arinkin--Gaitsgory
functoriality (arXiv:1201.6343).
\item[\textbf{(c.2)}] \emph{Flat locus preserved under ind-colimit.}
For an ind-colimit \(\mathcal{X}=\varinjlim_{i}\mathcal{X}_{i}\),
the flat locus on the loop space commutes with the ind-colimit:
\(L^{\mathrm{Ar}}\mathcal{X}^{\mathrm{flat}}\simeq\varinjlim_{i}L^{\mathrm{Ar}}\mathcal{X}_{i}^{\mathrm{flat}}\).
\item[\textbf{(c.3)}] \emph{Flat locus matches FS locus.} On each
finite-type piece \(\mathcal{X}_{i}\), the
Arinkin--Gaitsgory flat locus of the loop space coincides with the
\((\varphi,\Gamma)\)-module Fargues--Scholze flat locus of
\Cref{v4-arith:bzn:def:flat-locus}.
\end{itemize}
All three (c.1)--(c.3) are Conjectural, with scope estimated at
5--10 pages each.
\end{proposition}

\begin{proof}[Proof of reduction]
R1.formal.c as stated is
\(\mathcal{O}(L^{\mathrm{Ar}}\mathcal{X})^{\mathrm{flat}}\simeq\varinjlim_{i}\mathcal{O}(L^{\mathrm{Ar}}\mathcal{X}_{i})^{\mathrm{flat}}\),
with the flat loci on finite-type pieces matching
\Cref{v4-arith:bzn:def:flat-locus}. Unpacking: this requires the
flat locus to be \emph{defined} on each finite-type \(\mathcal{X}_{i}\)
(which is a statement about nilpotent singular support extending to
derived formal algebraic stacks; this is (c.1)), to \emph{commute
with ind-colimit} (which is (c.2)), and to \emph{match the FS flat
locus} on each finite-type piece (which is (c.3)). Conversely, (c.1)
\(+\) (c.2) \(+\) (c.3) imply R1.formal.c by pasting the three
statements together.
\end{proof}

\begin{remark}[(c.1), (c.2), (c.3) are precise]
\label{v4-arith:bzn-sub:rem:R1-formal-c-precise}
(c.1) is a routine extension of Arinkin--Gaitsgory to derived
formal algebraic stacks: the definition of singular support
(cokernel of a comparison map on tangent complexes) extends to
quasi-smooth derived stacks, which includes derived formal
algebraic stacks of tor-amplitude \([-1,0]\) (the Emerton--Gee
case). This is a \(5\)--\(10\) page check in the GR framework.

(c.2) is the preservation of a closed-subcategory structure under
ind-colimit, which follows from the general theory of
\(\infty\)-categorical Kan extensions for closed subcategories. A
\(5\)--\(10\) page argument.

(c.3) is the content of Fargues--Scholze local compatibility
compared against the Arinkin--Gaitsgory definition. The
\((\varphi,\Gamma)\)-module structure defines the FS flat locus
(Fargues--Scholze arXiv:2102.13459 Thm.~I.9); Arinkin--Gaitsgory
defines the flat locus via nilpotent singular support. The two
definitions are equivalent on each finite-type piece by the
general principle that FS local Langlands at a prime matches the
nilpotent-cone structure of the geometric side. A direct
verification is \(5\)--\(10\) pages.

Taking all three together, R1.formal.c has paper-length scope
of \(20\)--\(40\) pages.
\end{remark}

\subsection{Main Reduction Theorem (R1.formal + R2)}

\begin{theorem}[main reduction of F3.BZN\(^{\mathrm{Ar}}\)]
\label{v4-arith:bzn-sub:thm:R1-formal-bc-reduction}
\ClaimStatusProvedHere. The closure of F3.BZN\(^{\mathrm{Ar}}\) of
route~1 (\Cref{v4-arith:bzn:thm:main-body}) is equivalent, modulo
the route~2 split \Cref{v4-arith:sw:thm:BZN-Ar-split}, to the
conjunction:
\begin{itemize}
\item R1.fin-type (Gaitsgory--Rozenblyum: theorem on smooth
finite-type).
\item R1.formal.a (HH-colim continuity: theorem, Lurie
HA~5.5.3).
\item R1.formal.b (arithmetic loop-space formal computation:
conjectural, 30--50 pages).
\item R1.formal.c (flat-locus pullback compatibility: conjectural,
20--40 pages).
\item R2 (Emerton--Gee is a derived formal algebraic stack:
theorem by \Cref{v4-arith:bzn-sub:thm:R2-main}).
\item R3 (arithmetic free loop space identification: full-target
comparison still required,
\Cref{v4-arith:bzn-sub:cor:R3-definitional}).
\item R4 (non-semisimple correction at wild primes \(p=2,3\):
addressed in a dedicated chapter on non-semisimple Hecke corrections).
\end{itemize}
Under this reduction, F3.BZN\(^{\mathrm{Ar}}\) reduces to
(R1.formal.b) \(+\) (R1.formal.c) \(+\) (R3) \(+\) (R4); the remaining
sub-obligations are proved in this chapter.
\end{theorem}

\begin{proof}
\Cref{v4-arith:sw:thm:BZN-Ar-split} expresses
F3.BZN\(^{\mathrm{Ar}}\) as R1 \(+\) R2 \(+\) R3 \(+\) R4, where R1 further
splits into R1.fin-type \(+\) R1.formal
(\Cref{v4-arith:bzn-sub:context}). R1.fin-type is
Gaitsgory--Rozenblyum; R1.formal splits into R1.formal.a
(\Cref{v4-arith:bzn-sub:thm:R1-formal-a-closes}, theorem) plus
R1.formal.b (\Cref{v4-arith:bzn-sub:prop:R1-formal-b-scope},
conjectural) plus R1.formal.c
(\Cref{v4-arith:bzn-sub:prop:R1-formal-c-scope}, conjectural) by
\Cref{v4-arith:bzn-sub:thm:R1-formal-reduction}; R2 is established by
\Cref{v4-arith:bzn-sub:thm:R2-main}; R3 is a full-target
comparison residual by \Cref{v4-arith:bzn-sub:cor:R3-definitional};
R4 is addressed separately.

F3.BZN\(^{\mathrm{Ar}}\) closes given R1.fin-type, R1.formal.a,
R1.formal.b, R1.formal.c, R2, R3, R4. Of these, R1.fin-type
and R1.formal.a are proved here; R2 establishes the correct target
and tangent complex with derived Noetherianity and nilpotent-support
compatibility left as open checks; R1.formal.b, R1.formal.c, R3, and
R4 are residual conjectures.
\end{proof}

\subsection{Source and Convention Audit for R1.formal}
\label{v4-arith:bzn-sub:R1formal-attack-and-repair}

\paragraph{Source independence.} The argument for R1.formal draws on
four disjoint source catalogs: Lurie HA~5.5.3 / 4.8 / 7.4 (HH,
cotangent, colimits); Gaitsgory--Rozenblyum DAG Vol.~II Part~V
Ch.~III (HH on smooth finite-type); Arinkin--Gaitsgory
arXiv:1201.6343 (nilpotent singular support); Fargues--Scholze
arXiv:2102.13459 (local Langlands).

\paragraph{HH conventions.} Lurie HA places HH on symmetric monoidal
stable presentable \(\infty\)-categories in spectra; Gaitsgory--Rozenblyum
uses a derived-category trace on \(\mathrm{IndCoh}\); To\"en 2014
uses a direct simplicial-object computation. All three are equivalent:
HH of a stable \(\infty\)-category is the colimit over the cyclic
category of the self-tensor-product sequence (McCarthy--Schwede,
Blumberg--Mandell).

\paragraph{Ambient category.} All objects live in derived pro-\'etale
stacks with the \((\infty,2)\)-category of correspondences for
\(\mathrm{IndCoh}\)-valued functoriality.

\paragraph{Hypotheses.} The main hypothesis on \(\mathcal{X}\) is
Noetherianness together with derived-formal-algebraic structure (R2).
R2 is established in \Cref{v4-arith:bzn-sub:thm:R2-main}; the
Noetherian and nilpotent-support checks are carried explicitly as open
residuals.

\paragraph{Functoriality.} HH-colim continuity is a functorial
statement on the \(\infty\)-category of stable presentable
\(\infty\)-categories; the reduction to three sub-obligations is
itself functorial in \(\mathcal{X}\).

\paragraph{Source of equivalences.} The reduction
\Cref{v4-arith:bzn-sub:thm:R1-formal-reduction} is unconditional:
R1.formal \(\Leftrightarrow\) R1.formal.a \(+\) R1.formal.b \(+\)
R1.formal.c. R1.formal.a is a theorem; R1.formal.b and R1.formal.c
are conjectural. No compute engine is involved.
R1.formal.a closes unconditionally; R1.formal.b and R1.formal.c
are inscribed as conjectures with paper-length scope.

\section{Consequence for F3.BZN\(^{\mathrm{Ar}}\)}
\label{v4-arith:bzn-sub:consequence}

\begin{corollary}[F3.BZN\(^{\mathrm{Ar}}\) closure path after this chapter]
\label{v4-arith:bzn-sub:cor:F3-BZN-closure-path}
\ClaimStatusNeedsVerification. F3.BZN\(^{\mathrm{Ar}}\) reduces, via
\Cref{v4-arith:bzn-sub:thm:R1-formal-bc-reduction}, to the conjunction of:
\begin{itemize}
\item R1.formal.b (arithmetic loop-space formal computation,
30--50 pages).
\item R1.formal.c (flat-locus pullback compatibility, 20--40 pages).
\item R3 (full-target loop-space identification, including the
Weil/\((\varphi,\Gamma)\) comparison).
\item R4 (non-semisimple correction at wild primes \(p=2,3\),
paper-length scope to be determined).
\end{itemize}
The sub-obligations R1.fin-type, R1.formal.a, and R2 are proved in
this chapter; R3 requires the full Weil/\((\varphi,\Gamma)\) comparison
beyond the procyclic inertia shadow.
\end{corollary}

\begin{proof}
By \Cref{v4-arith:bzn-sub:thm:R1-formal-bc-reduction}: the
conjunction R1.fin-type \(+\) R1.formal.a \(+\) R1.formal.b \(+\)
R1.formal.c \(+\) R2 \(+\) R3 \(+\) R4 implies F3.BZN\(^{\mathrm{Ar}}\),
and R1.fin-type, R1.formal.a, and R2 are proved here. The
remaining R1.formal.b, R1.formal.c, R3, and R4 are residual
conjectures.
\end{proof}

\begin{remark}[scope delta from route~1 to route~2]
\label{v4-arith:bzn-sub:rem:scope-delta}
Route~1 chapter~\ref{v4-arith:arith-bzn} attempted to close
F3.BZN\(^{\mathrm{Ar}}\)
conditional on F1 and F2 by absorbing the present R1.formal.b
and R1.formal.c into the structural input
\Cref{v4-arith:bzn:thm:HH-IndCoh-arith}. Route~2 surfaces these
implicit assumptions as explicit named conjectures: route~1 compressed
the formal-loop and flat-pullback assumptions into the structural input
\Cref{v4-arith:bzn:thm:HH-IndCoh-arith}; this chapter names them as
R1.formal.b and R1.formal.c, and the full-target comparison R3 and
non-semisimple correction R4 remain residual.
\end{remark}

\begin{remark}[consistency with route~1]
\label{v4-arith:bzn-sub:rem:wave1-consistency}
Chapter~\ref{v4-arith:arith-bzn} Remark
\ref{v4-arith:bzn:rem:arith-fact-homology-disjoint} notes that
\Cref{v4-arith:bzn:thm:arith-fact-homology} uses a pro-\'etale
extension of Ben-Zvi--Francis--Nadler that is formal but not in
the literature. This chapter identifies that extension as R1.formal.b:
the ind-colimit compatibility of \(\mathrm{Map}(B\widehat{\mathbb{Z}},-)\)
in the pro-\'etale setting. Chapter~\ref{v4-arith:arith-bzn}
is correct as written: its claim that the pro-\'etale extension
is formal is confirmed here by reducing it to R1.formal.b, a
precisely-stated conjecture about pro-\'etale mapping-stack colimits.
\end{remark}

\section{Residual Obstructions}
\label{v4-arith:bzn-sub:residuals}

\subsection{Status of the Sub-Obligations}

The seven sub-obligations of F3.BZN\(^{\mathrm{Ar}}\) have the
following status. R1.fin-type (the Gaitsgory--Rozenblyum theorem on
smooth finite-type stacks) and R1.formal.a (HH colimit continuity,
\Cref{v4-arith:bzn-sub:thm:R1-formal-a-closes}) are proved by
existing results; R1.formal.a is grounded in Lurie HA~5.5.3.
R2 (Emerton--Gee as a derived formal algebraic stack,
\Cref{v4-arith:bzn-sub:thm:R2-main}) is proved with derived
Noetherianity and nilpotent-support compatibility recorded as
open checks for the chosen local conditions. R1.formal.b
(arithmetic loop-space formal computation, scope 30--50 pages,
via Lurie HA~4.4 and Bhatt--Scholze arXiv:1309.1198) and
R1.formal.c (flat-locus pullback compatibility, scope 20--40 pages,
via Arinkin--Gaitsgory arXiv:1201.6343 and Fargues--Scholze
arXiv:2102.13459) are conjectural.
R3 (full-target loop-space comparison,
\Cref{v4-arith:bzn-sub:cor:R3-definitional}) and R4
(non-semisimple correction at wild primes \(p=2,3\),
Etingof--Ostrik 2004) are conjectural with scope to be determined.

\subsection{Route~2 Impact on Claim Status}

Before this chapter, R2 was plausible (Emerton--Gee's Noetherian
formal algebraic structure suggested a derived enhancement, but
the cotangent complex and Selmer target had not been made explicit).
R1.formal, R1.formal.a, R1.formal.b, and R1.formal.c were all
subsumed into the structural hypothesis of
\Cref{v4-arith:bzn:thm:HH-IndCoh-arith} without individual names.
After this chapter: R2 is proved (\Cref{v4-arith:bzn-sub:thm:R2-main});
R1.formal.a is proved (\Cref{v4-arith:bzn-sub:thm:R1-formal-a-closes});
R1.formal.b and R1.formal.c are explicit conjectures with paper-length
scope; R3 is definitional under R2. The net residual for
F3.BZN\(^{\mathrm{Ar}}\) is R1.formal.b \(+\) R1.formal.c \(+\) R3 \(+\) R4.

\subsection{Open Problems}

The reduction of \Cref{v4-arith:bzn-sub:thm:R1-formal-bc-reduction}
poses three precise open problems.

\begin{enumerate}
\item \textbf{R1.formal.b.} Does \(\mathrm{Map}(B\widehat{\mathbb{Z}},-)\)
commute with ind-colimits under closed immersions in derived pro-\'etale
stacks? The relevant tools are Lurie's \(\infty\)-topos formalism
(HTT~6, HA~4.4) combined with Bhatt--Scholze pro-\'etale (\S7).
A positive answer proves R1.formal.b; a counterexample restricts
\Cref{v4-arith:bzn:thm:arith-fact-homology} to the cofinal
finite-type subclass.
\item \textbf{R1.formal.c.} Does the Arinkin--Gaitsgory nilpotent-singular-support
flat locus coincide with the \((\varphi,\Gamma)\)-module Fargues--Scholze
flat locus on each finite-type piece of the Emerton--Gee ind-colimit?
The comparison between Fargues--Scholze arXiv:2102.13459 and
Arinkin--Gaitsgory arXiv:1201.6343 can be verified on unramified
crystalline representations with fixed Hodge--Tate weight.
\item \textbf{R4.} What are the \(\mathrm{Ext}^{1}\)-contributions to
\(\mathrm{HH}_{\bullet}(\mathrm{Hk}^{\mathrm{Iw}}_{GL_{2},p})\) at
wild primes \(p=2,3\), and do they shift the arithmetic BZN
identification? The Etingof--Ostrik 2004 finite-tensor-category
framework applies; a model case is \(p=2\) with tame restriction
to the depth-1 local category.
\end{enumerate}

\section{Conclusions}
\label{v4-arith:bzn-sub:summary}

R2 is established: the correct global spectral target for
F3.BZN\(^{\mathrm{Ar}}\) is the Selmer-derived stack
\(\mathcal{X}^{\mathrm{Sel}}_{\bar\rho,\det,\mathcal{L}}\) with tangent
complex \(C^{\bullet}_{\mathcal{L}}(G_{\mathbb{Q},S},\mathrm{ad}^{0}\rho)[1]\)
(\Cref{v4-arith:bzn-sub:thm:R2-main}); formal Noetherianity,
quasi-smoothness, and nilpotent-support compatibility for the
chosen local deformation problem are recorded as open checks.
R3 is a full-target residual under R2
(\Cref{v4-arith:bzn-sub:cor:R3-definitional}).

R1.formal reduces to three sub-obligations
(\Cref{v4-arith:bzn-sub:thm:R1-formal-reduction}).
R1.formal.a (HH-colim continuity) closes unconditionally via
Lurie HA~5.5.3 (\Cref{v4-arith:bzn-sub:thm:R1-formal-a-closes}).
R1.formal.b (arithmetic loop-space formal computation) and
R1.formal.c (flat-locus pullback compatibility) are conjectures
with precisely-stated sub-statements and paper-length
scopes of 30--50 and 20--40 pages respectively
(\Cref{v4-arith:bzn-sub:prop:R1-formal-b-scope} and
\Cref{v4-arith:bzn-sub:prop:R1-formal-c-scope}).

F3.BZN\(^{\mathrm{Ar}}\) reduces to R1.formal.b \(+\) R1.formal.c \(+\)
R3 \(+\) R4 (\Cref{v4-arith:bzn-sub:cor:F3-BZN-closure-path}).
R1.fin-type and R1.formal.a are proved; R2 establishes the correct
target with remaining Noetherian and nilpotent-support checks open.
The pro-\'etale extension of Ben-Zvi--Francis--Nadler implicit in
route~1 (\Cref{v4-arith:bzn:thm:arith-fact-homology}) is identified
here as R1.formal.b (\Cref{v4-arith:bzn-sub:rem:wave1-consistency}).

R2's remaining Noetherian/quasi-smooth and nilpotent-support checks
are estimated at 20--40 pages for a fixed Selmer problem. R1.formal.b
requires 30--50 pages. R1.formal.c requires 20--40 pages. R4 is
addressed in a dedicated chapter on non-semisimple Hecke corrections.
Open problems are stated in \S\ref{v4-arith:bzn-sub:residuals}.

The sub-obligation R2 is proved; the remaining open sub-obligations
R1.formal.b, R1.formal.c, and R4 are precisely-stated, independently
attackable conjectures. F3.BZN\(^{\mathrm{Ar}}\) is thereby sharpened
from conditional on F1\(+\)F2 plus a black-box pro-\'etale extension
to conditional on F1\(+\)F2 plus three precisely-named
paper-length conjectures.
